\section{A Taxonomy of Coding Agents}
\label{sec:taxonomy}

To organize the diverse landscape of coding agents, we propose a taxonomy with three dimensions: autonomy level, interaction scope, and integration mode.

\subsection{Dimension 1: Autonomy Level}

We identify four levels of autonomy:

\begin{description}
    \item[Level 1 -- Reactive] The system responds only to explicit user requests. It takes no initiative and makes no suggestions unless prompted.
    \item[Level 2 -- Proactive] The system can suggest actions or improvements, but the human must execute any changes.
    \item[Level 3 -- Semi-autonomous] The system can execute actions, but requires human approval for each step or significant action.
    \item[Level 4 -- Autonomous] The system can plan and execute multi-step tasks independently, seeking human input only for clarification or final review.
\end{description}

\subsection{Dimension 2: Interaction Scope}

The scope of code that a system can understand and modify:

\begin{description}
    \item[Single-line] Completions and suggestions at the cursor position.
    \item[Function/block] Generation or modification of complete functions or code blocks.
    \item[File] Operations spanning an entire file, including refactoring and restructuring.
    \item[Repository] Cross-file changes, feature implementation, and codebase-wide modifications.
\end{description}

\subsection{Dimension 3: Integration Mode}

How the system integrates with the developer's workflow:

\begin{description}
    \item[IDE extension] Plugins for existing editors (VS Code, JetBrains, Vim).
    \item[Chat interface] Web-based or standalone conversation interfaces.
    \item[CLI tool] Command-line interfaces for terminal-centric workflows.
    \item[Standalone platform] Complete development environments with integrated AI capabilities.
\end{description}

\subsection{System Categories}

Using this taxonomy, we identify four primary categories. Table~\ref{tab:taxonomy} summarizes how each category maps to our three dimensions, where L1--L4 refer to the autonomy levels defined above.

\begin{table}[h]
\centering
\caption{Taxonomy of coding agent categories}
\label{tab:taxonomy}
\begin{tabular}{lccc}
\toprule
\textbf{Category} & \textbf{Autonomy} & \textbf{Scope} & \textbf{Integration} \\
\midrule
Code Completion Tools & L1--L2 & Single-line--Block & IDE extension \\
Conversational Assistants & L2--L3 & Function--File & Chat interface \\
Autonomous Agents & L3--L4 & File--Repository & CLI, Standalone \\
Integrated Environments & L1--L4 & Single-line--Repository & Standalone platform \\
\bottomrule
\end{tabular}
\end{table}

The following sections examine each category in detail.
